\section{Radical placement and scale attachment}
\label{sec:radical-placement}
% Status: open. States the radical-placement problem and source-to-operator obligations.

The secular equation contains the radical
\begin{equation}
  \Delta_J^{1/2}=\sqrt{J^2+4J}.
  \label{eq:radical-object}
\end{equation}
The manuscript treats \(\Delta_J^{1/2}\) as the central object whose physical address must be identified.  The algebra alone specifies a two-branch spectrum.  Physics begins when the radical is assigned to a pole condition, a boundary determinant, a compactification eigenvalue, a worldsheet constraint, or a running matching condition.

\subsection{Placement options}

The possible placements are distinct research programs.
\begin{center}
\small
\begin{tabular}{p{0.20\linewidth}p{0.28\linewidth}p{0.42\linewidth}}
\toprule
Placement & Physical meaning & Required derivation\\
\midrule
Pole spectrum & \(\Delta_J^{1/2}\) enters a propagator or self-energy pole condition & derive the W/Z pole ratio and scalar partner branch from one spectral equation\\
GUT boundary & \(\Delta_J^{1/2}\) fixes a high-scale weak-angle datum & embed the \(J\)-samples in a unification representation and run to the pole scale\\
Running parameter & \(\Delta_J^{1/2}\) constrains a scheme-dependent coupling or vacuum parameter & specify the scheme, scale, matching map, and threshold corrections\\
KK eigenvalue & \(\Delta_J^{1/2}\) comes from a compact internal wave operator & identify the internal operator, mode labels, degeneracies, and boundary conditions\\
Open-string endpoint & \(\Delta_J^{1/2}\) comes from endpoint or Chan--Paton data & derive the two-channel matrix from endpoint charges and brane/bulk mixing\\
Worldsheet constraint & \(\Delta_J^{1/2}\) comes from level matching, conformal weights, or a modular sector & identify the CFT labels and show how they reduce to the electroweak determinant\\
\bottomrule
\end{tabular}
\end{center}

The pole-spectrum placement is the current working reading.  It gives the numerical clue its electroweak content because W and Z are observed as poles.  It also gives the negative branch a task: the same determinant should know about the scalar/order-parameter side.  This reading demands the most direct dynamical calculation, since it asks string/Kaluza--Klein structure to survive down to a low-energy pole observable.

The GUT-boundary placement gives a familiar comparison class.  In SU(5) and Spin(10) settings, representation theory can give a clean weak-angle value at the high scale, with running and thresholds supplying the low-energy relation \cite{Britto2021}.  A DeVries high-scale reading would join that lineage.  The current manuscript pursues the pole reading because the negative branch and pole-ratio clue appear together in the same low-energy normalization.

\subsection{Scale attachment}

Let \(\mathcal O_J\) denote the physical operator whose two-channel reduction gives
\begin{equation}
  \det(\lambda-\mathcal O_J)=\lambda^2+J\lambda-J.
  \label{eq:operator-placement}
\end{equation}
The scale question is then a question about where \(\mathcal O_J\) lives:
\begin{equation}
  \mathcal O_J
  \in
  \left\{
  \mathcal O_{\rm pole},
  \mathcal O_{\rm GUT},
  \mathcal O_{\rm KK},
  \mathcal O_{\rm brane},
  \mathcal O_{\rm CFT}
  \right\}.
  \label{eq:operator-addresses}
\end{equation}
Each address has a different interpretation of the same algebra.  A pole operator must be built from four-dimensional gauge-Higgs correlators.  A GUT operator must be built from representation traces and matching.  A KK operator must be a compact internal differential operator or boundary value problem.  A brane operator must couple localized and bulk modes.  A worldsheet operator must arise from conformal weights, current algebra, or level matching.

This list fixes the scheme dependence of the radical.  A different placement changes the required derivation.

\subsection{Placement chain}
\label{sec:placement-chain}

The placement question can be recorded as a chain.  A source theory \(\mathcal S\) supplies an operator, the operator supplies a two-channel kernel, the kernel supplies branches, and the branches acquire physical meaning after a scheme and an assignment rule have been fixed:
\begin{equation}
\begin{aligned}
\mathcal S
&\longrightarrow
\mathcal O_J
\longrightarrow
\mathcal M_J
\longrightarrow
\{x_+(J),x_-(J)\},
\\
\{x_+(J),x_-(J)\}
&\longrightarrow
\left\{
\sPole,\,
\mathcal F_{\rm sc},\,
\mathcal C_{\rm top}
\right\}.
\end{aligned}
\label{eq:placement-chain}
\end{equation}
Here \(\mathcal F_{\rm sc}\) denotes the scalar/order-parameter functional sought in Sec.~\ref{sec:negative}, and \(\mathcal C_{\rm top}\) denotes the global-form and representation constraints discussed in Secs.~\ref{sec:flavour} and \ref{sec:global-form}.  The chain separates four tasks:
\begin{align}
  \text{operator origin:}\quad&
  \mathcal S\mapsto \mathcal O_J,
  \label{eq:chain-origin}\\
  \text{two-channel reduction:}\quad&
  \mathcal O_J\mapsto
  \begin{pmatrix}0&\sqrt J\\ \sqrt J&-J\end{pmatrix},
  \label{eq:chain-reduction}\\
  \text{branch reading:}\quad&
  \mathcal M_J\mapsto \{x_+(J),x_-(J)\},
  \label{eq:chain-branch}\\
  \text{physical assignment:}\quad&
  \{x_\pm\}\mapsto
  \left(\sPole,\mathcal F_{\rm sc},\mathcal C_{\rm top}\right).
  \label{eq:chain-assignment}
\end{align}

Each arrow has different source support.  Complex-pole and scheme language belongs to the last arrow \cite{MartinRobertson2025}.  Higgs doublet, adjoint current, custodial structure, and the electroweak ray belong to the assignment arrow \cite{Logan2022,Dawson2018}.  Endpoint labels, D-brane gauge fields, and interval boundary conditions belong to the first two arrows \cite{TongString,Polchinski1994,CsakiHubiszMeade2005}.  Equation~\eqref{eq:chain-branch} is algebraically controlled; the other arrows are open derivation targets.

This chain also organizes the negative branch.  If a candidate source theory derives \(\mathcal M_J\), then \(x_-(J)\) is carried by the same light-sector basis as \(x_+(J)\).  The scalar question becomes a question about the eigenvector map,
\begin{equation}
  u_-(J)\in\mathcal H_J
  \quad\longmapsto\quad
  \mathcal F_{\rm sc},
  \label{eq:negative-eigenvector-chain}
\end{equation}
where \(\mathcal H_J\) is the two-dimensional light sector.  A Higgs interpretation requires \(\mathcal F_{\rm sc}\) to be gauge invariant.  A brane interpretation requires \(\mathcal F_{\rm sc}\) to be a modulus, separation, or localized boundary datum.  A compactification interpretation requires \(\mathcal F_{\rm sc}\) to be a normalized internal eigenmode or deformation parameter.

The chain gives a compact criterion for acceptance.  A proposed mechanism should fill one of the arrows in Eq.~\eqref{eq:placement-chain} or state a new arrow with its source theory, scheme, and physical assignment.

\subsection{Original research route}

The route developed here is:
\begin{equation}
\begin{aligned}
  \text{Regge/high-spin datum}
  &\longrightarrow
  \text{two-channel secular determinant}
  \\
  &\longrightarrow
  \text{electroweak pole ratio and scalar partner branch}.
\end{aligned}
  \label{eq:research-route}
\end{equation}
The first arrow belongs to the string/dual-model tradition: spectra organize into towers with controlled asymptotics.  The second arrow is the DeVries algebra.  The third arrow is the open electroweak calculation.

A concrete model should therefore supply four ingredients:
\begin{enumerate}
\item an internal label whose quadratic invariant is \(J\);
\item a two-channel system whose off-diagonal mixing is fixed by \(\sqrt J\);
\item a diagonal trace term giving the \(-J\) entry;
\item a physical normalization that places the positive branch on the vector pole ratio and the negative branch on the scalar/order-parameter side.
\end{enumerate}

The fourth ingredient is the new conceptual burden.  It turns the algebraic clue into a physical program and fixes the relation between the pole reading and high-scale weak-angle boundary data.
